\documentclass[12pt,a4paper]{article}

\usepackage[T2A]{fontenc}
\usepackage[utf8]{inputenc}
\usepackage[serbianc]{babel} % Ćirilica
\usepackage{graphicx}
\usepackage{xcolor}
\usepackage{amsmath, amsthm}
\usepackage{wrapfig} % Za sliku sa leve strane

% Definisanje teoreme na srpskom
\newtheorem{teorema}{Теорема}

\title{Кратак преглед вештачке интелигенције}
\author{Андреј Шобић}
\date{\today}

\begin{document}

\maketitle

\tableofcontents
\newpage

\section{Увод}
Вештачка интелигенција (АИ) представља једну од најузбудљивијих области савремене компјутерске науке. Циљ је стварање система који могу да решавају проблеме који су се раније сматрали искључиво људском доменом.

\section{Математичке основе}
Вештачка интелигенција се у великој мери ослања на математику, посебно на вероватноћу и линеарну алгебру.

\begin{teorema}[Основни принцип учења]
Ако је систем добро подешен, грешка приликом предвиђања тежи нули како број итерација расте.
\end{teorema}

Једна од најпознатијих једначина у физици која се често користи као пример за једноставне моделе јесте позната Ајнштајнова формула:
$$E = mc^2$$

А у машинском учењу често минимизујемо функцију трошка, где је један од облика:
$$J(\theta) = \frac{1}{2m} \sum_{i=1}^{m} (h_\theta(x^{(i)}) - y^{(i)})^2$$

\section{Примена и форматирање}

\begin{wrapfigure}{l}{0.3\textwidth}
    \centering
    \includegraphics[width=0.25\textwidth]{slike/profil.jpg}
    \caption{Симболична слика}
\end{wrapfigure}

Текст може бити написан у различитим форматима. На пример, можемо писати \textbf{подебљано}, затим \textit{искошено}, или пак користити \textcolor{blue}{плаву боју} за наглашавање одређених делова текста.

Као што је приказано са леве стране, слика се може лепо уклопити у сам пасус без прављења непотребног празног простора.

\subsection{Врсте учења}
Постоје различити формати листа. Ево једне набрајачке листе:
\begin{itemize}
    \item Надгледано учење
    \item Ненадгледано учење
\end{itemize}

А ево и нумерисане листе за редослед корака:
\begin{enumerate}
    \item Прикупљање података
    \item Тренирање модела
\end{enumerate}

\subsection{Табеларни приказ}
У следећој табели је приказан кратак преглед:

\begin{center}
\begin{tabular}{|c|c|}
\hline
\textbf{Алгоритам} & \textbf{Тачност} \\ \hline
Линеарна регресија & 85\% \\ \hline
Неуронске мреже & 98\% \\ \hline
\end{tabular}
\end{center}

\section{Закључак}
Вештачка интелигенција се развија невероватном брзином. Разумевање њених основа је кључно за сваког будућег инжењера информатике.

\end{document}